\chapter{Politiques de groupe (GPOs)}

\section{GPO-PasswordPolicy}

\textbf{Appliquée sur :} Racine du domaine

\subsection{Politique de mot de passe}

\begin{table}[H]
    \centering
    \begin{tabularx}{\textwidth}{|p{3cm}|l|p{4cm}|}
        \hline
        \textbf{Paramètre} & \textbf{Valeur} & \textbf{Justification} \\
        \hline
        Enforce password history & 24 & Empêche réutilisation immédiate \\
        \hline
        Max password age & 90 jours & Force changement régulier \\
        \hline
        Min password age & 45 jours & Empêche changement trop fréquent \\
        \hline
        Min password length & 10 chars & Force sécurité suffisante \\
        \hline
        Complexity required & Enabled & Maj, min, chiffre, symbole \\
        \hline
        Reversible encryption & Disabled & Mots de passe restent hashés \\
        \hline
    \end{tabularx}
    \caption{Politique de mot de passe}
    \label{tab:pwd-policy}
\end{table}

\subsection{Politique de verrouillage de compte}

\begin{table}[H]
    \centering
    \begin{tabularx}{\textwidth}{|l|l|}
        \hline
        \textbf{Paramètre} & \textbf{Valeur} \\
        \hline
        Lockout threshold & 5 tentatives \\
        \hline
        Lockout duration & 30 minutes \\
        \hline
        Lockout counter reset & 30 minutes \\
        \hline
    \end{tabularx}
    \caption{Politique de verrouillage}
    \label{tab:lockout-policy}
\end{table}

\section{GPO-BaseSecurity}

\textbf{Appliquée sur :} Tous les postes

Paramètres de sécurité de base :

\begin{itemize}
    \item \textbf{Windows Defender Firewall} : Activé sur tous les profils (Domain, Private, Public)
    \item \textbf{Windows Defender} : Forcé actif (désactiver l'option "Turn off Defender")
    \item \textbf{SMBv1} : Désactivé (Network → Lanman Workstation)
    \item \textbf{AutoRun USB} : Désactivé (AutoPlay Policies → Turn off AutoPlay)
    \item \textbf{Screen saver} : Timeout 10 minutes + mot de passe
\end{itemize}

\section{GPO-Bitlocker}

\textbf{Appliquée sur :} Postes sensibles (RH, Comptabilité)

\begin{lstlisting}
Computer Configuration
  → Policies
    → Administrative Templates
      → Windows Components
        → Encryption
          → BitLocker Drive Encryption
\end{lstlisting}

\begin{table}[H]
    \centering
    \begin{tabularx}{\textwidth}{|p{3cm}|p{4cm}|}
        \hline
        \textbf{Paramètre} & \textbf{Valeur} \\
        \hline
        Encryption method & AES 256 \\
        \hline
        Drive encryption type & Used space only \\
        \hline
        Startup auth & TPM + PIN (si pas TPM) \\
        \hline
        Recovery key storage & Active Directory \\
        \hline
    \end{tabularx}
    \caption{Configuration BitLocker}
    \label{tab:bitlocker}
\end{table}

\section{GPO-Direction}

Héritée de BaseSecurity + PasswordPolicy

\begin{itemize}
    \item Internet via proxy (filtrage SquidGuard)
    \item Réseaux sociaux : autorisé
    \item Streaming vidéo : bloqué
    \item Téléchargements : bloqués
\end{itemize}

\section{GPO-Commercial}

Héritée de BaseSecurity + PasswordPolicy

\begin{itemize}
    \item Internet via proxy (filtrage SquidGuard)
    \item Réseaux sociaux : autorisé
    \item Email access : autorisé
    \item Téléchargements : bloqués
    \item VLAN 20 (isolation)
\end{itemize}

\section{GPO-RH}

Héritée de BaseSecurity + PasswordPolicy + BitLocker

\subsection{Blocage USB}

\begin{lstlisting}
Computer Configuration > Policies > Windows Settings > Security
  → System Services → USBSTOR

Action       : Update
Hive         : HKEY_LOCAL_MACHINE
Key Path     : SYSTEM\CurrentControlSet\Services\USBSTOR
Value name   : start
Value type   : REG_DWORD
Value data   : 4 (disabled)
\end{lstlisting}

\begin{itemize}
    \item Internet restrictif (email uniquement)
    \item Réseaux sociaux : bloqués
    \item Téléchargements : bloqués
    \item BitLocker obligatoire
    \item USB bloqué (USBSTOR désactivé)
    \item VLAN 50 (isolation)
\end{itemize}

\section{GPO-Comptabilite}

Héritée de BaseSecurity + PasswordPolicy + BitLocker

\begin{itemize}
    \item Internet HTTPS uniquement (pas HTTP)
    \item Tous les téléchargements bloqués
    \item Réseaux sociaux : bloqués
    \item BitLocker obligatoire
    \item USB bloqué (USBSTOR désactivé)
    \item VLAN 60 (isolation maximale)
\end{itemize}

\section{GPO-Informatique}

\textbf{Block Inheritance : OUI} (ignore les politiques parentes)

\begin{itemize}
    \item Accès panel de contrôle complet
    \item Gestionnaire tâches accessible
    \item WSUS obligatoire (mises à jour depuis srv-dc01:8530)
    \item Pas de proxy (accès Internet direct)
    \item Accès inter-VLAN complet (administration)
\end{itemize}

\section{GPO-Serveurs}

Appliquée sur OU=Serveurs (DCs, NAS)

\begin{itemize}
    \item Print Spooler : \textbf{désactivé} (DCs n'impriment pas)
    \item Firewall Windows : activé
    \item Advanced audit logging
    \item Event log retention : 7 jours minimum
\end{itemize}

\section{PSO (Password Settings Objects)}

Permet plusieurs politiques dans le même domaine.

\begin{table}[H]
    \centering
    \begin{tabularx}{\textwidth}{|l|l|l|l|}
        \hline
        \textbf{Groupe} & \textbf{Min Length} & \textbf{Max Age} & \textbf{Lockout} \\
        \hline
        IT Admins & 15 & 60 j & 3 \\
        \hline
        Utilisateurs & 10 & 90 j & 5 \\
        \hline
        Temp Services & 12 & 30 j & 5 \\
        \hline
    \end{tabularx}
    \caption{PSO Configuration}
    \label{tab:pso}
\end{table}

\section{Vérification et application}

\begin{lstlisting}[style=powershell]
# Forcer update GPO
gpupdate /force

# Voir GPOs appliquées
gpresult /h report.html

# Tester GPO sur machine
Get-GPResultantSetOfPolicy -ComputerName PC-BUR01
\end{lstlisting}
